\section{Tensor Schwinger functions}\label{sec:tensor-axioms}

The Gaussian estimate of \cite[Proposition~6.1]{ContinuumLimits} defines a continuous map
$\Psi\mapsto Q_\Psi$ from $\Sspace_n$ to $L^2(\lambda)$.
To use law-level reflection positivity for arbitrary tensor tests,
we must show that support in a half-space makes this random polynomial
measurable with respect to the corresponding half-space sigma algebra. The Hermite
expansion used for the growth bound does not preserve supports.
We therefore use a compact Fourier expansion for this step.

\subsection{Products and supported tensor tests}

For $\Psi\in\Sspace_k$ and $\Phi\in\Sspace_l$, the tensor evaluation
obeys
\begin{equation}\label{eq:tensor-product-rule}
 Q_\Psi Q_\Phi=Q_{\Psi\otimes\Phi}\quad\text{in }L^1(\lambda).
\end{equation}
Both sides are separately continuous bilinear maps to $L^1$:
for the left side use the $L^2$ bounds and Cauchy--Schwarz, and for
the right side use continuity of the tensor product in Schwartz
seminorms. They agree for elementary tensors. Their linear span
is dense, so approximation in each variable proves
\eqref{eq:tensor-product-rule}. The same argument gives
\begin{equation}\label{eq:tensor-equivariance}
 Q_{\overline\Psi}=\overline{Q_\Psi},\qquad
 Q_{g\cdot\Psi}=\widehat g Q_\Psi
 \quad\text{if }g_*\lambda=\lambda.
\end{equation}
Here the second identity uses the exterior-square action, including
the tensor indices.

We recall the needed supported expansion. Let $O\subset\R^3$ be open
and $\Psi\in C_c^\infty(O^n)$. Its support projects into compact
sets $K_a\subset O$. Choose $\chi_a\in C_c^\infty(O)$ equal to one
on $K_a$, and a cube of side $s$ containing all these supports in
its interior. Extend $\Psi$ periodically from this cube in all its
variables. Its Fourier coefficients $c_{\mathbf k}$ obey, for
every integer $p$,
\[
 |c_{\mathbf k}|\le C_{p,s,\Psi}(1+|\mathbf k|^2)^{-p},
 \qquad \mathbf k=(k_1,\ldots,k_n)\in(\Z^3)^n.
\]
This follows by applying $(1-\Delta)^p$ and integrating by parts;
there are no boundary terms. Multiplication of the Fourier series by
$\prod_a\chi_a(y_a)$ gives the identity on all space
\begin{equation}\label{eq:supported-tensor-expansion}
 \Psi=\sum_{\mathbf k}c_{\mathbf k}
       \bigotimes_{a=1}^n
       \left(\chi_a(y_a)e^{2\pi i k_a\cdot y_a/s}\right).
\end{equation}
The seminorm $q_m$ of each tensor factor product grows at most as
$(1+|\mathbf k|^2)^{m/2}$. Taking $p=m+3n$ makes the series
absolutely convergent in that seminorm. The same argument applies
componentwise to two-form tensor tests. It also proves the density
of elementary tensors used above; arbitrary Schwartz tests are
approximated by growing compact cutoffs.

Each elementary tensor in \eqref{eq:supported-tensor-expansion}
has factors supported in $O$, so its random evaluation is
$\Sigma_O$-measurable and belongs to $L^2(\lambda)$ by the moment
bound. Continuity of $Q$ and closedness of $L^2(\Sigma_O,\lambda)$
give
\begin{equation}\label{eq:tensor-local-measurability}
 Q_\Psi\in L^2(\Sigma_O,\lambda)
       \qquad(\Psi\in C_c^\infty(O^n)).
\end{equation}

This extends to tests supported in the closed half-space
$H=\{x:v\cdot x\ge c\}$, with $|v|=1$. Choose a fixed compact
cutoff $\chi$ equal to one near zero and put
\[
 \Psi_j(y)=\prod_{a=1}^n\chi(y_a/j)\,
           \Psi(y_1-v/j,\ldots,y_n-v/j).
\]
Each $\Psi_j$ has compact support in
$O^n$, $O=\{x:v\cdot x>c\}$. Translation is continuous in Schwartz
space, and the growing cutoff tends to one in its action on each
Schwartz function. More explicitly, Leibniz' rule gives
\[
 q_m(\Psi_j-\Psi)\le C_{m,n,\chi}j^{-1}q_{m+1}(\Psi).
\]
Consequently $Q_{\Psi_j}\to Q_\Psi$ in $L^2(\lambda)$, and
\begin{equation}\label{eq:tensor-closed-halfspace}
 Q_\Psi\in L^2(\Sigma_O,\lambda)\qquad
                   (\supp\Psi\subset H^n).
\end{equation}
Thus tests meeting the boundary plane belong to the same
positive-side $L^2$ space.

\subsection{The five conditions}

\begin{proposition}\label{prop:tensor-OS}
Let $\lambda$ be a Gaussian-dominated law of real tempered two-forms.
Suppose it is invariant under Euclidean motions and reflection
positive at every plane through the origin. Then its tensor
Schwinger functions satisfy \textup{(E0$'$)}--\textup{(E3)} of
\cref{thm:main}. If, for every $A,B\in L^2(\lambda)$,
\[
 \E_\lambda[\overline A\,U(a)B]\longrightarrow
       \overline{\E_\lambda A}\,\E_\lambda B
       \quad\text{as }|a|\longrightarrow\infty,
\]
they also satisfy \textup{(E4)}.
\end{proposition}

\begin{proof}
The Hermite estimate of \cite[Proposition~6.1]{ContinuumLimits} proves \textup{(E0$'$)}. Covariance follows from
\eqref{eq:tensor-equivariance} and invariance of expectation.
Permutation symmetry holds on elementary tensors because the random
field evaluations commute; continuity gives \textup{(E3)} on all
tensor tests.

For \textup{(E2)}, let $r$ be the stated reflection and $H$ the
chosen closed half-space. By \eqref{eq:tensor-closed-halfspace},
$A=\sum_{k=0}^MQ_{\Psi_k}$ lies in its positive-side $L^2$ space,
where $Q_{\Psi_0}=\Psi_0$. Reflection positivity extends from
bounded characters to this whole space by their density and the
boundedness of the reflection form. Equations
\eqref{eq:tensor-product-rule}--\eqref{eq:tensor-equivariance} give
\[
 \overline{\mathcal R_rQ_{\Psi_k}}
       =Q_{\Theta_k\Psi_k},\qquad
 Q_{\Theta_k\Psi_k}Q_{\Psi_l}
       =Q_{\Theta_k\Psi_k\otimes\Psi_l}\quad\text{in }L^1.
\]
Therefore
\[
 0\le\E_\lambda[(\overline{\mathcal R_rA})A]
   =\sum_{k,l=0}^MS_{k+l}(\Theta_k\Psi_k\otimes\Psi_l).
\]
This is the tensor reflection form with the conjugation and signs
in \cref{thm:main}.

Finally, \eqref{eq:tensor-product-rule} and translation covariance
identify the expression in \textup{(E4)} with
$\E_\lambda[Q_\Psi U(a)Q_\Phi]$. Both random polynomials are in
$L^2(\lambda)$. Apply the stated clustering with
$A=\overline{Q_\Psi}$ and $B=Q_\Phi$ to obtain
$\E_\lambda Q_\Psi\,\E_\lambda Q_\Phi$.
\end{proof}

For the components in \cref{prop:limit-nine-mirrors}, the clustering
hypothesis holds. The covariance theorem in \cref{sec:rotations}
will supply their Euclidean invariance and reflection positivity
at every plane.


